\documentclass[reqno,11pt]{amsart}
\usepackage{math327-lecture}
\usepackage{needspace}
\usepackage{tikz}
\newcommand{\SO}{\operatorname{SO}}
\lectureinfo{24}{Nov. 4, 2026}{Rotations, Euler's Theorem, and Poles of Finite Rotation Groups}
\artinsections{5.1; 6.12}

% Typeset from handwritten Lecture 17, Oct. 10, 2025, pp. 1--10,
% and handwritten Lecture 18, Oct. 15, 2025, pp. 5--16.
\begin{document}
\raggedbottom
\makelecturetitle



\section{Rotations and the special orthogonal group}

Recall that
\[
  \SO_3(\R)=\{A\in\GL_3(\R):A^{\mathsf T}A=I,\ \det A=1\}.
\]
This is a group under matrix multiplication.

\begin{lemma}[Preservation of dot products]\label{lem:dot-product}
If $A\in\SO_3(\R)$ and $v,w\in\R^3$, then
$(Av)\cdot(Aw)=v\cdot w$.
\end{lemma}
\begin{proof}
Writing $v\cdot w=v^{\mathsf T}w$, we have
\[
  (Av)\cdot(Aw)=v^{\mathsf T}A^{\mathsf T}Aw
  =v^{\mathsf T}w=v\cdot w.
\]
\end{proof}

\begin{define}[Rotation and pole]
A \emph{rotation} of $\R^3$ is a linear map $\rho:\R^3\to\R^3$
that fixes a unit vector $u$ and acts as a planar rotation on
$u^\perp$. The vector $u$ is called a \emph{pole} of $\rho$;
the line $\R u$ is its axis.
\end{define}

\Needspace{11\baselineskip}
\begin{lemma}\label{lem:fixed-vector}
Every $A\in\SO_3(\R)$ has $1$ as an eigenvalue.
\end{lemma}
\begin{proof}
Since $\det A=1$ and $AA^{\mathsf T}=I$,
\begin{align*}
  \det(A-I)
  &=\det(A^{\mathsf T}-I)\\
  &=\det\bigl(A(A^{\mathsf T}-I)\bigr)\\
  &=\det(I-A)=-\det(A-I).
\end{align*}
The last sign occurs because the matrices are $3\times3$.
Thus $\det(A-I)=0$, so $Au=u$ for some nonzero $u$.
\end{proof}

\begin{theorem}[Euler's rotation theorem]\label{thm:euler}
A linear map $\rho:\R^3\to\R^3$ is a rotation if and only if its
matrix in the standard basis belongs to $\SO_3(\R)$.
\end{theorem}
\begin{proof}
Suppose first that $\rho$ is a rotation with pole $u$.
Choose an orthonormal basis $\mathcal B=(u,v,w)$, where $(v,w)$
is an orthonormal basis of $u^\perp$. In this basis its matrix is
\[
  B_\theta=
  \begin{pmatrix}
    1&0&0\\
    0&\cos\theta&-\sin\theta\\
    0&\sin\theta&\cos\theta
  \end{pmatrix}.
\]
Let $P$ have columns $u,v,w$. Since $P$ is orthogonal, the matrix
in the standard basis is $PB_\theta P^{-1}=PB_\theta P^{\mathsf T}$.
This matrix is orthogonal and has determinant $1$.

Conversely, let $A\in\SO_3(\R)$. By Lemma~\ref{lem:fixed-vector},
there is a unit vector $u$ with $Au=u$. Put $W=u^\perp$.
For $w\in W$, Lemma~\ref{lem:dot-product} gives
\[
  (Aw)\cdot u=(Aw)\cdot(Au)=w\cdot u=0,
\]
so $A$ preserves $W$. In an orthonormal basis adapted to
$\R^3=\R u\oplus W$, its matrix has the form
\[
  \begin{pmatrix}1&0\\0&B\end{pmatrix},
\]
where $B$ is a $2\times2$ orthogonal matrix and $\det B=\det A=1$.
An orthogonal matrix in dimension $2$ with determinant $1$ has the
form
\[
  B=\begin{pmatrix}\cos\theta&-\sin\theta\\
                       \sin\theta&\cos\theta\end{pmatrix}.
\]
Hence $A$ fixes $u$ and rotates its perpendicular plane.
\end{proof}

\begin{remark}
Orthogonal maps of a plane have determinant $1$ or $-1$.
The first are rotations and the second are reflections.
The determinant condition above is what rules out a reflection
on $W$.
\end{remark}

\section{Poles of a finite rotation group}

Let $G\leq\SO_3(\R)$ be finite, and write $N=|G|$.
We want to study $G$ by letting it act on the poles of its rotations.

\begin{define}[Poles of a group]
The set of \emph{poles of $G$} is
\[
  \mathcal P=\{p\in S^2:gp=p\text{ for some }g\in G\setminus\{1\}\},
\]
where $S^2=\{p\in\R^3:\|p\|=1\}$ is the unit sphere.
\end{define}

A nonidentity rotation fixes exactly two points on the unit sphere,
the two unit vectors along its axis. Indeed, its restriction to the
perpendicular plane is a nonidentity planar rotation and fixes no
nonzero vector. Since $G$ is finite, $\mathcal P$ is finite as well.

\Needspace{18\baselineskip}
\begin{example}[Poles of the tetrahedral group]
Let $T$ be the group of rotational symmetries of a regular
tetrahedron centered at the origin. Its poles are the directions
from the origin through the following points:

\begin{minipage}[c]{0.53\linewidth}
\centering
\begin{tabular}{lr}
  \textbf{Direction through} & \textbf{Poles}\\[3pt]
  A vertex & $4$\\
  A face center & $4$\\
  An edge midpoint & $6$\\[2pt]
  \hline
  Total & $14$
\end{tabular}
\end{minipage}\hfill
\begin{minipage}[c]{0.43\linewidth}
\centering
\begin{tikzpicture}[scale=0.8]
  \draw[gray] (0,0) circle (1.65);
  \draw[dashed,gray!65] (-1.65,0) arc (180:360:1.65 and 0.35);
  \draw[gray!65] (1.65,0) arc (0:180:1.65 and 0.35);
  \coordinate (a) at (0,1.429);
  \coordinate (b) at (1.556,-0.476);
  \coordinate (c) at (-0.778,0.197);
  \coordinate (d) at (-0.778,-1.150);
  \draw[dashed,courseblue] (c)--(a) (c)--(b) (c)--(d);
  \draw[thick,courseblue] (a)--(b)--(d)--cycle;
  \foreach \v in {a,b,c,d} {\fill[courseblue] (\v) circle (2pt);}
\end{tikzpicture}
\end{minipage}

The axes through a vertex and the opposite face center give
rotations of order $3$. The axes through midpoints of opposite
edges give half-turns. A pole is the point on the unit sphere in
the indicated direction, rather than the face center or edge
midpoint itself.
\end{example}

\begin{proposition}\label{prop:action-on-poles}
The natural action of $G$ on $S^2$ restricts to an action on
$\mathcal P$.
\end{proposition}
\begin{proof}
Let $p\in\mathcal P$, and choose $g\neq1$ with $gp=p$.
For $h\in G$, the nonidentity element $hgh^{-1}\in G$ satisfies
\[
  (hgh^{-1})(hp)=hgp=hp.
\]
Also, $hp$ is a unit vector because $h$ preserves lengths.
Thus $hp\in\mathcal P$.
\end{proof}

\section{Counting poles and stabilizers}

Decompose the pole set into distinct orbits:
\[
  \mathcal P=\mathcal O_{p_1}\sqcup\cdots\sqcup\mathcal O_{p_k}.
\]
Write
\[
  n_i=|\mathcal O_{p_i}|,\qquad r_i=|G_{p_i}|.
\]
Every $r_i\geq2$, since $p_i$ is fixed by a nonidentity element.
By orbit--stabilizer,
\[
  n_i r_i=N.
\]
The stabilizers of points in the same orbit are conjugate, so their
orders agree.

\begin{theorem}[Pole-counting formula]\label{thm:pole-counting}
For a finite subgroup $G\leq\SO_3(\R)$ of order $N$,
\[
  \boxed{\displaystyle
  \sum_{p\in\mathcal P}\bigl(|G_p|-1\bigr)=2(N-1).}
\]
\end{theorem}
\begin{proof}
Count the pairs $(g,p)$ with $g\in G\setminus\{1\}$,
$p\in\mathcal P$, and $gp=p$.
For each $p$, there are $|G_p|-1$ choices of $g$.
For each nonidentity $g$, there are exactly two choices of $p$,
its two poles. The two counts give the formula.
\end{proof}

Grouping the sum by orbits gives the equivalent identities
\[
  \sum_{i=1}^k n_i(r_i-1)=2(N-1),
  \qquad
  \sum_{i=1}^k\left(1-\frac1{r_i}\right)=2-\frac2N.
\]
These combine the orbit sizes and stabilizer orders into a
restriction on finite rotation groups, the starting point for
their classification.

\section{Consequences of the pole-counting formula}

Suppose throughout that $G$ is nontrivial. Since $r_i\geq2$,
each term $1-1/r_i$ is at least $1/2$ and strictly less than $1$.
The right side of the pole-counting formula lies in $[1,2)$.
Consequently the poles form either two or three $G$-orbits:
one orbit would contribute less than $1$, while four would
contribute at least $2$.

\begin{lemma}[Point stabilizers are cyclic]\label{lem:cyclic-stabilizer}
For every $p\in\mathcal P$, the stabilizer $G_p$ is cyclic.
\end{lemma}
\begin{proof}
Every element of $G_p$ fixes $p$ and restricts to an
orientation-preserving orthogonal map of $p^\perp$. Restriction
embeds $G_p$ into $\SO_2(\R)$, the group of planar rotations.
A finite subgroup of $\SO_2(\R)$ is cyclic: choose its smallest
positive rotation angle $\theta$; division with remainder by
$\theta$ shows that every angle is an integral multiple of it.
\end{proof}

\subsection*{Two orbits}
If $k=2$, then
\[
  \frac1{r_1}+\frac1{r_2}=\frac2N.
\]
Since $r_i\leq N$, each summand on the left is at least $1/N$.
Equality forces $r_1=r_2=N$, hence $n_1=n_2=1$.
The pole set is therefore $\{p,-p\}$, and every element of $G$
fixes both poles. Thus $G=G_p$ is cyclic by
Lemma~\ref{lem:cyclic-stabilizer}. Geometrically, these are the
rotations about one common axis.

\subsection*{Three orbits}
For $k=3$, order the stabilizer sizes as $2\leq r_1\leq r_2\leq r_3$.
The formula becomes
\[
  \frac1{r_1}+\frac1{r_2}+\frac1{r_3}=1+\frac2N>1.
\]
This forces $r_1=2$: three reciprocals with denominators at
least $3$ sum to at most $1$. It also forces $r_2\leq3$,
since $1/2+1/4+1/4=1$. If $r_2=2$, write $r_3=m\geq2$.
If $r_2=3$, then $r_3<6$, so $r_3$ is $3$, $4$, or $5$.
Using $n_i=N/r_i$ gives all the numerical possibilities:
\[
\begin{array}{c|c|c}
 (r_1,r_2,r_3) & N=|G| & (n_1,n_2,n_3)\\ \hline
 (2,2,m),\ m\geq2 & 2m & (m,m,2)\\
 (2,3,3) & 12 & (6,4,4)\\
 (2,3,4) & 24 & (12,8,6)\\
 (2,3,5) & 60 & (30,20,12)
\end{array}
\]
The first row is the dihedral pattern. The remaining rows are
realized by rotational symmetry groups of the tetrahedron,
cube or octahedron, and dodecahedron or icosahedron,
respectively. The pole-counting argument gives the possible
orders and orbit sizes; identifying a group from these data
requires a geometric argument.

\begin{example}[The tetrahedral pattern]
For the tetrahedral rotation group $T$ from the earlier example,
the six edge-midpoint poles have stabilizers of order $2$.
The four vertex poles and four face-center poles each have
stabilizers of order $3$. Thus the three orbit sizes are
$(6,4,4)$ and $|T|=12$, as in the $(2,3,3)$ row.

Let $V$ be one of the four-point orbits, and choose $z\in V$.
The order-$3$ stabilizer $T_z$ cyclically permutes the other
three points of $V$. Indeed, a nonidentity rotation fixing $z$
can fix on the sphere only $z$ and $-z$. If $-z$ lay in $V$,
the remaining two points would also have to be fixed by an
order-$3$ permutation, a contradiction. Thus $T_z$ acts as a
$3$-cycle on $V\setminus\{z\}$, and $T$ is transitive on ordered
pairs of distinct points of $V$.
Since rotations preserve distances, every pair of distinct
points of $V$ has the same distance: the four points are the
vertices of a regular tetrahedron inscribed in $S^2$.
\end{example}

\end{document}
